\begin{center}
\textsc{Gauss} an \textsc{Olbers}.\quad Göttingen, 14.~Januar 1827.
\end{center}

\par\noindent\hspace*{\parindent}\mbox{.\ .\ .\ .\ .\ .}\quad Ich glaube, Ihnen schon früher mitgeteilt zu haben, dass ich es für ganz undurchführbar halte, das Werk, das ich künftig über meine Messungen veröffentlichen möchte, auch in theoretischer Hinsicht ganz eigenständig zu machen, wenn ich nicht wenigstens einen großen Teil der Theorie zuvor anderswo eigens behandle.
Schon die Darlegung der Gründe für mein Vorgehen wird viel Raum einnehmen, aber, wenn ich mich so ausdrücken darf, die \textls[120]{Gründe} der \textls[120]{Gründe} können nicht in das Werk selbst aufgenommen werden, ohne dass es ganz bunt zusammengewürfelt und doch unbefriedigend wird.
Ich habe mich deshalb entschlossen, verschiedene theoretische Themen zunächst in einzelnen Abhandlungen gesondert zu entwickeln. Nur so ist es auch möglich, diese bedeutenden neuen Kapitel der Mathematik mit einem gewissen Grad an Vollständigkeit auszuarbeiten.
Gewissermaßen ist meine Preisschrift über die Transformation der Flächen die erste dieser Abhandlungen. Die zweite habe ich vor einigen Monaten der k.\ Sozietät als Supplementum theoriae combinationis observationum etc. übergeben.
Sie enthält die Prinzipien, die der \textls[120]{Ausgleichung} der Beobachtungen zugrunde gelegt werden müssen, und sogar einige Beispiele aus den Messungen \textsc{Krayenhoffs} und meinen eigenen.\newline
\mbox{.\ .\ .\ .\ .\ .}\par
